\chapter*{General Introduction}
\addcontentsline{toc}{chapter}{General Introduction}

In a world where digital transformation accelerates at an unprecedented rhythm, cybersecurity has imposed itself as a \textbf{major strategic imperative} for enterprises and institutions alike. The explosion of cyberattacks --- particularly sophisticated ransomware campaigns and massive data breaches, which cost an average of \$4.45 million per incident according to IBM --- reveals a \textit{systemic threat} that reaches far beyond the purely technical sphere. These challenges thrust organisations into a critical predicament: strengthening their security posture is no longer an option; it is a vital necessity to ensure long-term sustainability, preserve reputation, and comply with ever more demanding regulations.

Faced with this urgency, the \textbf{Security Operations Center (SOC)} constitutes the beating heart of cyber defence. A true operational control tower, it delivers real-time proactive threat detection, in-depth incident analysis, and coordinated response. Yet building a traditional SOC represents a \textit{formidable challenge} for the majority of organisations --- especially small and medium-sized enterprises (SMEs) and resource-constrained institutions. Exorbitant infrastructure costs, software licence fees, recruitment of scarce experts, the complexity of managing heterogeneous technologies, and a global shortage of cybersecurity talents (3.4 million unfilled positions according to ISC\textsuperscript{2}) render this model either inaccessible or ineffective for many.

It is precisely against this backdrop that an \textit{innovative paradigm} emerges: \textbf{SOC as a Service (SOCaaS)}. This disruptive model outsources the core functions of a conventional SOC --- continuous monitoring, threat intelligence, incident analysis, and response --- through a secure, cloud-native, fully managed offering.

Driven by these strategic imperatives, our final-year project is anchored within \textit{Algérie Telecom (AT)}with a clear operational ambition: to design and deploy an open-source SOCaaS that combines \textbf{industrial performance}, \textbf{economic accessibility}, and \textbf{technological sovereignty}. This initiative answers an undeniable observation: local SMEs and institutions, despite being vital to the digital ecosystem, remain vulnerable to increasingly sophisticated cyber threats simply because they lack dedicated security resources. Our solution positions itself as a credible alternative to proprietary offerings --- offerings that are often ill-adapted to tight budgets and the specific needs of these stakeholders.

The proposed architecture transcends mere technical provision; it embodies a \textit{lever for cultural transformation}. By relying exclusively on open-source building blocks --- \textbf{Wazuh} for SIEM and XDR capabilities, \textbf{TheHive} and \textbf{Shuffle} for SOAR orchestration --- we guarantee not only total transparency of security mechanisms but also lasting vendor independence. This approach fosters genuine appropriation of security tools by local teams, transforming every user into a potential contributor to the system's continuous improvement --- a virtuous cycle that serves collective resilience.

Beyond detection and incident response, our SOCaaS integrates a structured educational component: training paths adapted to system administrators, playbooks contextualised to the beneficiaries' core business, and awareness workshops focused on targeted risks such as ransomware and phishing. This \textbf{human dimension} is the true keystone of our vision: to cultivate a mature cyber culture where security becomes a shared reflex --- not an externally imposed constraint.

This thesis articulates the approach across four chapters:

\begin{itemize}
    \item \textbf{Chapter 1: Project Framework} --- describes the professional context in which this project is inscribed.
    \item \textbf{Chapter 2: Overview of the Security Operations Center} --- explores the fundamentals of SOCs, existing service models, and associated technological solutions.
    \item \textbf{Chapter 3: Design of the SOCaaS Solution} --- dedicated to technology selection, architecture modelling, and integration of the open-source stack.
    \item \textbf{Chapter 4: Realization, Deployment, and Validation} --- details the infrastructure deployment, functional tests, and the results obtained under realistic operational conditions.
\end{itemize}